% AY2024 力学 解答例
% 体裁・粒度は 参考/ のPDFに揃える（行間を埋め、最終解答は \ans{} で boxed）。
% 解答・数値・問題は本年度過去問から自力で導いた（東大版は体裁の手本のみ）。
\documentclass[11pt,a4paper]{article}
\usepackage{../../../00_common/preamble}

\usetikzlibrary{decorations.pathmorphing}

\title{力学 AY2024 解答例（専門基礎科目）}
\author{}
\date{}

\begin{document}
\maketitle

% ==================================================================
\section*{第1問〔I〕}

ばね定数 $k\,(k>0)$ のばねの一端を壁に固定し，他端に質量 $m$ の物体を取り付け，
なめらかな水平床上に置く。ばね自然長のときの物体重心位置を $x=0$ とする。
自然長から $x$ 正向きに $L\,(L>0)$ だけ伸ばして静かに手を離すと，物体は速度に
比例する空気抵抗（比例定数 $B>0$）を受けて $x$ 軸上を運動する。

\subsection*{(1) 運動方程式}
物体には次の二つの水平力がはたらく。
\begin{itemize}
  \item ばねの復元力：伸び $x$ に比例し，変位と逆向きなので $-kx$。
  \item 空気抵抗：速度 $\dot{x}$ に比例し，運動を妨げる向きなので $-B\dot{x}$。
\end{itemize}
床はなめらかなので摩擦はない。ニュートンの運動方程式 $m\ddot{x}=\sum F$ より
\[
  m\ddot{x} = -kx - B\dot{x}.
\]
\[
  \ans{\; m\ddot{x} + B\dot{x} + kx = 0 \;}
\]
初期条件は $x(0)=L,\ \dot{x}(0)=0$ である。

\subsection*{(2) 往復運動せず原点に近づく条件}
(1) は定数係数の 2 階線形同次微分方程式である。解を $x=e^{\lambda t}$ と仮定して
代入すると，特性方程式
\begin{align*}
  m\lambda^2 + B\lambda + k = 0
  \quad\therefore\quad
  \lambda = \frac{-B \pm \sqrt{B^2 - 4mk}}{2m}
\end{align*}
を得る。物体が「往復運動せず（振動せず）時間と共に原点に近づく」ためには，
解が三角関数（振動）を含まず，かつ $t\to\infty$ で $x\to0$ となればよい。

\medskip
\noindent\textbf{(i) 振動しない条件}：判別式が負だと
$\lambda$ は虚部をもち $e^{\lambda t}$ に $\cos,\sin$ が現れて往復運動する。
したがって振動しないためには根が実数，すなわち
\[
  B^2 - 4mk \ge 0.
\]

\noindent\textbf{(ii) 原点に近づく条件}：このとき二根 $\lambda_1,\lambda_2$ はともに実数で，
\[
  \lambda_1\lambda_2 = \frac{k}{m}>0,\qquad
  \lambda_1+\lambda_2 = -\frac{B}{m}<0
\]
より両根とも負である（$B>0,k>0$）。ゆえに $x(t)$ は単調に減衰し，
$t\to\infty$ で $x\to0$ となる。したがって条件は (i) だけでよい。

\medskip
$B,m,k>0$ より $B^2\ge 4mk \iff B \ge 2\sqrt{mk}$。
\[
  \ans{\; B^2 \ge 4mk \quad\Bigl(\text{すなわち } B \ge 2\sqrt{mk}\Bigr) \;}
\]
（等号 $B=2\sqrt{mk}$ が臨界減衰，$B>2\sqrt{mk}$ が過減衰である。）

\medskip
検算: 無次元の減衰比 $\zeta=\dfrac{B}{2\sqrt{mk}}$ を用いると，振動しない条件は
$\zeta\ge1$ と書ける。これは $B\ge2\sqrt{mk}$ と一致する。
また次元を確認すると，$[m k]=\mathrm{kg}\cdot(\mathrm{kg/s^2})=\mathrm{kg^2/s^2}$，
$\sqrt{mk}$ は $\mathrm{kg/s}$ で抵抗係数 $B$（$\mathrm{N\cdot s/m=kg/s}$）と同次元であり整合する。✓

% ==================================================================
\clearpage
\section*{第2問〔II〕}

半径 $r$，慣性モーメント $I$ の滑車を天井に固定し，糸を巻き付けて一端に質量 $m$ の
物体を取り付ける。静止状態から静かに手を離すと，物体は鉛直下向きに落下する。
重力加速度の大きさを $g$，糸の質量・太さは無視する。糸は滑車上で滑らないとする。

\subsection*{(1) 物体の加速度}
鉛直下向きを正にとり，物体の加速度を $a$，糸の張力を $T$，滑車の角加速度を
$\alpha$ とする。糸が滑らないので拘束条件は
\[
  a = r\alpha .
\]
物体の並進の運動方程式（下向き正）：
\[
  m a = mg - T. \tag{1}
\]
滑車の回転の運動方程式（張力 $T$ が腕 $r$ で滑車を回す）：
\[
  I\alpha = T r
  \quad\therefore\quad
  T = \frac{I\alpha}{r} = \frac{I a}{r^2}. \tag{2}
\]
(2) を (1) に代入して
\begin{align*}
  m a = mg - \frac{I a}{r^2}
  \;\Longrightarrow\;
  \left(m + \frac{I}{r^2}\right) a = mg
  \;\Longrightarrow\;
  a = \frac{mg}{\,m + I/r^2\,}.
\end{align*}
\[
  \ans{\; a = \dfrac{m g r^2}{\,m r^2 + I\,} \;}
\]

\subsection*{(2) 糸の張力}
(2) より $T=\dfrac{I a}{r^2}$ に (1) で求めた $a$ を代入すると
\begin{align*}
  T = \frac{I}{r^2}\cdot\frac{m g r^2}{m r^2 + I}
    = \frac{m g I}{\,m r^2 + I\,}.
\end{align*}
\[
  \ans{\; T = \dfrac{m g I}{\,m r^2 + I\,} \;}
\]

\subsection*{(3) 落下距離 $h$ での滑車の角速度}
物体は等加速度 $a$ で落下するので，初速 $0$，落下距離 $h$ のとき速さ $v$ は
\[
  v^2 = 2 a h.
\]
糸が滑らないので物体の速さと滑車外周の速さは等しく $v=r\omega$，すなわち
$\omega=v/r$。よって
\begin{align*}
  \omega^2 = \frac{v^2}{r^2} = \frac{2 a h}{r^2}
           = \frac{2h}{r^2}\cdot\frac{m g r^2}{m r^2 + I}
           = \frac{2 m g h}{\,m r^2 + I\,}.
\end{align*}
\[
  \ans{\; \omega = \sqrt{\dfrac{2 m g h}{\,m r^2 + I\,}} \;}
\]

\medskip
検算: エネルギー保存で確認する。物体が $h$ 落下したとき
\begin{align*}
  m g h = \tfrac{1}{2}m v^2 + \tfrac{1}{2}I\omega^2
        = \tfrac{1}{2}m (r\omega)^2 + \tfrac{1}{2}I\omega^2
        = \tfrac{1}{2}(m r^2 + I)\omega^2 .
\end{align*}
これを解くと $\omega^2=\dfrac{2mgh}{mr^2+I}$ となり，(3) の結果と一致する。✓
また $I\to0$（質量・慣性のない滑車）の極限では $a\to g$，$T\to0$，
$\omega\to\sqrt{2gh}/r$（自由落下 $v=\sqrt{2gh}$）となり物理的に妥当。✓

\end{document}
